\documentclass[11pt]{scrartcl}
\usepackage[a4paper,margin=2.6cm]{geometry}
\usepackage{../common/rlstudy}

\title{Extended reading --- Week 18: RLHF \& PPO for Language Models}
\date{\today}
\begin{document}
\maketitle

\section*{How to read this list}
Each paper below gets: full citation, one-paragraph summary, a
\emph{Why it matters for this week} note, and one annotated key
passage with brief commentary. The list is ordered chronologically.

\section*{1. Schulman et al.\ 2017 --- Proximal Policy Optimization}
\textbf{Citation.} J. Schulman, F. Wolski, P. Dhariwal, A. Radford,
O. Klimov, \emph{Proximal Policy Optimization Algorithms},
arXiv:1707.06347, 2017.

\paragraph{Summary.} Introduces the clipped surrogate
$\mathcal{L}^{\text{CLIP}}$ and the adaptive-KL variant
$\mathcal{L}^{\text{KLPEN}}$ as first-order, off-the-shelf substitutes
for TRPO. Demonstrates competitive performance on continuous control
(MuJoCo) and Atari with a tiny implementation footprint. Becomes the
default policy-gradient algorithm in deep RL within twelve months of
release.

\paragraph{Why it matters for this week.} Every line of our PPO
implementation is in this paper. The clipped surrogate, the
GAE-$\lambda$ advantages, and the adaptive KL controller all originate
here. Even modern RLHF variants like PPOv2 in TRL are conservative
re-wrappings of this 2017 algorithm.

\paragraph{Annotated key passage.}
\begin{quote}\itshape
``By taking the minimum of the clipped and unclipped objective, the
final objective is a lower bound (i.e., a pessimistic bound) on the
unclipped objective.''
\end{quote}
\noindent This is the single sentence that justifies Theorem~3.1
(``Why the clip is a pessimistic surrogate'') in the notes; it also
explains why PPO is robust to large advantages without exploding.

\section*{2. Ouyang et al.\ 2022 --- InstructGPT (the foundational RLHF paper)}
\textbf{Citation.} L. Ouyang \emph{et al.}, \emph{Training language
models to follow instructions with human feedback},
arXiv:2203.02155, NeurIPS 2022.

\paragraph{Summary.} The InstructGPT paper. Trains a 1.3B-parameter
GPT-3 variant in three stages --- SFT on demonstrations, Bradley--Terry
reward modelling on human preferences, PPO against the reward model
with a KL penalty against the SFT reference --- and shows it is
preferred by human raters over 175B GPT-3 on a broad range of
instruction-following tasks. Forms the template for ChatGPT, Claude,
Gemini and essentially every subsequent production aligned LLM.

\paragraph{Why it matters for this week.} The three-stage pipeline
(the three-stage pipeline in \texttt{notes.pdf}, Algorithm~1) is theirs. Their Figure~5 (preference
win-rate vs.\ model size) is the empirical justification for the
entire field. They are also the first to publicly enumerate the
specific failure modes --- sycophancy, refusal collapse, length bias ---
that we'll spend weeks 18--22 mitigating.

\paragraph{Annotated key passage.}
\begin{quote}\itshape
``We find that the outputs of our 1.3B InstructGPT model are preferred
to those of the 175B GPT-3, despite having 100x fewer parameters.''
\end{quote}
\noindent The headline result. The implication --- that an
alignment-trained small model can dominate a much larger base model
on instruction-following --- shifts the whole research economics of
language modelling toward post-training.

\section*{3. Gao, Schulman \& Hilton 2023 --- Scaling laws for reward overoptimization}
\textbf{Citation.} L. Gao, J. Schulman, J. Hilton, \emph{Scaling Laws
for Reward Model Overoptimization}, PMLR v202 (ICML 2023).

\paragraph{Summary.} Constructs a controlled environment in which a
``gold'' reward model is treated as ground truth and a smaller
``proxy'' RM is fit on samples from it. PPO is then run against the
proxy while gold reward is measured on held-out data. Finds an
analytic functional form:
$R_{\text{gold}}(d) \approx d (\alpha - \beta d)$ for KL distance
$d \!=\! \sqrt{\KL(\pi_\theta\|\pi_{\text{ref}})}$, with both $\alpha$ and the
optimal $d^\star$ scaling smoothly with proxy-RM size and dataset
size.

\paragraph{Why it matters for this week.} This is \emph{the} paper on
reward hacking. It quantifies precisely why a KL leash and early
stopping are not optional. The schematic figure in
the notes is a stylised version of their Figure~1.

\paragraph{Annotated key passage.}
\begin{quote}\itshape
``We find that the gold reward, $R(d)$, can be very accurately fit by
a function of the form $R(d) = d(\alpha_{\text{bon}} -
\beta_{\text{bon}}\, d)$ for best-of-$n$ sampling, and $R(d) =
d(\alpha_{\text{RL}} - \beta_{\text{RL}}\, \log d)$ for reinforcement
learning.''
\end{quote}
\noindent The two functional forms differ --- RL is logarithmic in
$d$, BoN linear --- and the RL form decays \emph{faster} once KL is
past its sweet spot. This explains why brute-force PPO without
gold-eval is dangerous in production: you cannot see the gold drop
without independent evaluations.

\section*{4. Bai et al.\ 2022 --- Constitutional AI and RLAIF}
\textbf{Citation.} Y. Bai \emph{et al.\ }(Anthropic), \emph{Constitutional
AI: Harmlessness from AI Feedback}, arXiv:2212.08073, 2022.

\paragraph{Summary.} Replaces the human-preference half of RLHF with
\emph{AI feedback}. Two stages: (i) the model critiques and revises
its own outputs against a written ``constitution''; (ii) a separate
preference model is trained on AI-generated comparisons and used to
drive the RL loop, exactly as a human-RM would. The result, Claude's
harmlessness profile, is the first credible evidence that scalable
oversight via AI feedback can replace human feedback for
non-edge-case alignment.

\paragraph{Why it matters for this week.} Same machinery as
InstructGPT --- PPO, KL leash, BT loss --- but a completely different
source of supervision. Sets the stage for week~19's iterative DPO
and week~20's weak-to-strong supervision.

\paragraph{Annotated key passage.}
\begin{quote}\itshape
``[CAI] reduces the reliance on human supervision while preserving the
ability to fine-tune for harmlessness, suggesting a path to
scalable oversight of AI systems.''
\end{quote}
\noindent The plausibility claim of the entire scalable-oversight
agenda. The 2024--2026 literature (process reward models, deliberative
alignment, debate, etc.) all live downstream of this paper.

\section*{5. Lambert et al.\ 2024 --- T\"ulu 3}
\textbf{Citation.} N. Lambert \emph{et al.}, \emph{T\"ulu 3: Pushing
Frontiers in Open Language Model Post-Training}, arXiv:2411.15124,
2024.

\paragraph{Summary.} The first fully open RLHF+RLVR recipe at
production scale: a complete SFT $\to$ DPO $\to$ RLVR (PPO with
verifiable rewards) pipeline for Llama-3.1, with open datasets, open
code, and an open evaluation suite. Closes the gap with proprietary
models on a wide range of benchmarks while making every artefact
reproducible.

\paragraph{Why it matters for this week.} The closest thing in 2026
to a ``how to run RLHF on a real model, from scratch, end-to-end''
reference. The dataset (\texttt{tulu-3-sft-mixture},
\texttt{tulu-3-pref-mixture}) is what we'd reach for on a serious
RLHF experiment. The TRL recipe in
\texttt{project/trl\_recipe.py} is a scaled-down version of theirs.

\paragraph{Annotated key passage.}
\begin{quote}\itshape
``We provide a recipe that achieves competitive performance with
frontier closed models such as GPT-4o-mini and Claude 3.5 Haiku
across a wide range of benchmarks, while fully releasing our data,
model weights, training and evaluation code.''
\end{quote}
\noindent Open-source RLHF caught up with closed-source RLHF in
2024. The lesson is that the algorithmic core has stabilised --- the
remaining differentiator is data quality, not method.

\newpage
\section*{6. Christiano et al.\ 2017 --- Deep RL from Human Preferences}
\textbf{Citation.} P. Christiano, J. Leike, T. Brown, M. Martic,
S. Legg, D. Amodei, \emph{Deep Reinforcement Learning from Human
Preferences}, arXiv:1706.03741, NeurIPS 2017.

\paragraph{Summary.} The original RLHF paper. Trains an RL agent on
Atari and MuJoCo tasks using only pairwise human preferences over
trajectory segments --- no hand-engineered reward function. Lays
out the three-stage pipeline (preference collection, RM training,
RL against the RM) that would become the InstructGPT template five
years later.

\paragraph{Why it matters for this week.} The pipeline architecture
we use today was invented \emph{here}, in 2017, on Atari, by the
same OpenAI--DeepMind team that would later build ChatGPT. Reading
this paper alongside InstructGPT is the cleanest way to see the
continuity of ideas across the deep-RL-to-LLM transition.

\paragraph{Annotated key passage.}
\begin{quote}\itshape
``We show that this approach can effectively solve complex RL tasks
without access to the reward function, using less than 1\% of the
feedback that the agent receives.''
\end{quote}
\noindent The astonishing data efficiency point: a few hundred
human comparisons could replace millions of environment-reward
samples. The same calculus is what makes preference RL economical
for LLMs.

\section*{7. Ziegler et al.\ 2019 --- Fine-tuning language models from human preferences}
\textbf{Citation.} D.\,M.\ Ziegler, N. Stiennon, J. Wu, T.\,B.\
Brown, A. Radford, D. Amodei, P. Christiano, G. Irving,
\emph{Fine-Tuning Language Models from Human Preferences},
arXiv:1909.08593, 2019.

\paragraph{Summary.} The first paper to apply RLHF to a language
model (GPT-2). Pioneered the practical machinery --- KL-shaped
per-token rewards, value-head sharing the policy trunk, adaptive
KL --- that would later go into InstructGPT.

\paragraph{Why it matters for this week.} This is the paper that
turned RLHF from an RL idea (Christiano 2017) into an LLM
methodology. Almost every implementation detail you'll meet in TRL
or OpenRLHF was invented here.

\paragraph{Annotated key passage.}
\begin{quote}\itshape
``We use a KL constraint to keep the policy from drifting too far
from the supervised initialisation; without it, the policy quickly
collapses to a degenerate solution that the reward model gives high
score to but that humans rate as poor.''
\end{quote}
\noindent A clear early statement of reward-model overoptimisation,
written before Gao et al.\ 2023 quantified it.

\section*{8. Stiennon et al.\ 2020 --- Learning to summarize from human feedback}
\textbf{Citation.} N. Stiennon, L. Ouyang, J. Wu, D. Ziegler, R.
Lowe, C. Voss, A. Radford, D. Amodei, P. Christiano, \emph{Learning
to Summarize from Human Feedback}, arXiv:2009.01325, NeurIPS 2020.

\paragraph{Summary.} The TL;DR-summarisation benchmark for RLHF.
Trains a 6.7B model with PPO against a learned reward model on
human comparisons of news-article summaries. The RLHF-trained model
is preferred to a 100$\times$ larger supervised baseline.

\paragraph{Why it matters for this week.} The bridge between
\textsc{Ziegler 2019} (proof of concept) and \textsc{Ouyang 2022}
(industrial scale). Also the source of the much-cited
``reward-model overoptimisation'' plot, where held-out human
ratings degrade as RM-score keeps rising.

\paragraph{Annotated key passage.}
\begin{quote}\itshape
``As training progresses, the policy's reward according to the
reward model continues to improve, but its quality according to
human labels saturates and eventually degrades. This is a clear
indication of reward-model overoptimisation and motivates the
KL constraint in our algorithm.''
\end{quote}
\noindent The smoking-gun figure for the entire reward-hacking
literature. Read this paper carefully before deciding how to
length-regulate your own runs.

\newpage
\section*{Pointers for further depth}

\begin{itemize}[leftmargin=*]
  \item \textsc{Schulman et al.\ 2016} --- \emph{Generalized
        Advantage Estimation}, arXiv:1506.02438. The GAE
        bias-variance dial; if you want to understand the
        $\lambda$ knob in your RLHF runs, read this.
  \item \textsc{Casper et al.\ 2023} --- \emph{Open Problems and
        Fundamental Limitations of RLHF}, arXiv:2307.15217. A
        clear-eyed survey of why RLHF is much harder than it looks.
  \item \textsc{Coste et al.\ 2023} --- \emph{Reward Model Ensembles
        Help Mitigate Overoptimization}, arXiv:2310.02743. Ensembles
        of RMs vote; the policy is rewarded only when they agree.
  \item \textsc{Razin et al.\ 2024} --- \emph{Likelihood Displacement},
        arXiv:2410.08847. The technical explanation of why DPO and
        PPO both decrease $\log\pi(y_l)$ aggressively.
  \item \textsc{Lambert} --- \emph{The RLHF Book} (online, 2026). The
        best free reference text on the entire pipeline, regularly
        updated.
\end{itemize}

\newpage
\section*{Sourced from}
\texttt{research/llm\_rl\_sota.md} §1 (Classical RLHF), §6
(Constitutional AI / RLAIF), §7 (reward hacking, KL control), and
§11 (key reading list).

\end{document}
